\chapter{格林函数}
\label{ch:green}

\begin{flushright}
\textit{“举一隅不以三隅反，则不复也。”——《论语·述而》}
\footnote{格林函数是那一隅,即单位点源在给定边界下的响应.知道它以后,线性叠加把任意源的场一并还回来.}
\end{flushright}

分离变量把定解问题拆成本征模态,再逐个定系数.\term{格林函数}{gelinhanshu}走另一条路:先求单位点源在同样边界和初值条件下的场,再由线性叠加得到任意源的场.点源的影响函数称为格林函数.经典文献见~\cite{green1828,duffy,stakgold,jackson}.

这一章的泊松方程写成
\begin{equation}
  \nabla^2 u=f.
  \label{eq:gf-poisson-sign}
\end{equation}
因此自由空间的基本解是负的,
\begin{equation}
  \nabla^2\Bigl(-\frac{1}{4\pi R}\Bigr)=\delta^3(\br).
\end{equation}
若把方程改写成 $-\nabla^2 u=\rho$,格林函数整体差一个负号,下面的公式要一并改号.

\section{泊松方程的格林函数}
\idx{格林恒等式}{gelinhengdengshi}

\subsection{格林恒等式}
设 $u,v$ 在区域 $T$ 上有连续的二阶导数,边界 $\Sigma$ 分段光滑.对 $u\nabla v$ 用散度定理,
\begin{equation}
  \iint_{\Sigma} u\,\nabla v\cdot\mathrm{d}\mathbf{S}
  =\iiint_{T}\bigl(u\nabla^2 v+\nabla u\cdot\nabla v\bigr)\,\mathrm{d}V.
\end{equation}
这是格林第一恒等式.把 $u$ 和 $v$ 对调再相减,梯度内积消去,得到格林第二恒等式
\begin{equation}
  \iint_{\Sigma}\Bigl(
    u\frac{\partial v}{\partial n}-v\frac{\partial u}{\partial n}
  \Bigr)\,\mathrm{d}S
  =\iiint_{T}\bigl(u\nabla^2 v-v\nabla^2 u\bigr)\,\mathrm{d}V.
  \label{eq:green2}
\end{equation}
$\partial/\partial n$ 是外法向导数.

\subsection{点源与挖去奇点}
泊松方程 $\nabla^2 u=f$,边界上给定
\begin{equation}
  \Bigl[\alpha\frac{\partial u}{\partial n}+\beta u\Bigr]_{\Sigma}=\varphi.
  \label{eq:gf-bc}
\end{equation}
$\alpha=0$ 是第一类(狄利克雷),$\beta=0$ 是第二类(诺伊曼),两者都不为零是第三类(罗宾).

单位点源的场 $v(\br,\br_0)$ 满足 $\nabla^2 v=\delta(\br-\br_0)$.它在 $\br_0$ 处不可导,不能直接把 $u$ 和 $v$ 代入\eqref{eq:green2}.挖去以 $\br_0$ 为心、半径 $\varepsilon$ 的小球 $K_{\varepsilon}$.在剩下的区域上 $v$ 调和,$\nabla^2 u=f$,于是
\begin{equation}
\begin{aligned}
  &\iint_{\Sigma}\Bigl(v\frac{\partial u}{\partial n}-u\frac{\partial v}{\partial n}\Bigr)\,\mathrm{d}S
  +\iint_{\Sigma_{\varepsilon}}\Bigl(v\frac{\partial u}{\partial n}-u\frac{\partial v}{\partial n}\Bigr)\,\mathrm{d}S \\
  &\qquad=\iiint_{T\setminus K_{\varepsilon}} v f\,\mathrm{d}V.
\end{aligned}
\end{equation}
小球附近 $v\sim -1/(4\pi|\br-\br_0|)$.外法向在 $\Sigma_{\varepsilon}$ 上指向球心,$\partial/\partial n=-\partial/\partial\varepsilon$.面积是 $4\pi\varepsilon^2$,所以
\begin{equation}
  \iint_{\Sigma_{\varepsilon}} v\frac{\partial u}{\partial n}\,\mathrm{d}S\to 0,
  \qquad
  \iint_{\Sigma_{\varepsilon}} u\frac{\partial v}{\partial n}\,\mathrm{d}S\to -u(\br_0).
\end{equation}
第二式的来源是 $\partial v/\partial\varepsilon\sim 1/(4\pi\varepsilon^2)$,再乘面积 $4\pi\varepsilon^2$ 得到 $1$,法向反向贡献一个负号.令 $\varepsilon\to 0$,
\begin{equation}
\begin{aligned}
  u(\br_0)
  &=\iiint_{T} v(\br,\br_0)\,f(\br)\,\mathrm{d}V \\
  &\quad-\iint_{\Sigma}\Bigl[
      v\frac{\partial u}{\partial n}-u\frac{\partial v}{\partial n}
    \Bigr]\,\mathrm{d}S.
\end{aligned}
\label{eq:gf-fundamental}
\end{equation}
这个公式表面上同时要用边界上的 $u$ 和 $\partial u/\partial n$.定解问题只给它们的一个线性组合.把 $v$ 选成满足同样类型的齐次边界条件,未知的那一组边界值就消掉.这样选出的 $v$ 记作格林函数 $G$.

\subsection{三类边界}
狄利克雷问题要求 $G|_{\Sigma}=0$.\eqref{eq:gf-fundamental}变成
\begin{equation}
  u(\br_0)
  =\iiint_{T} G(\br,\br_0)\,f(\br)\,\mathrm{d}V
  +\iint_{\Sigma}\phi(\br)\frac{\partial G}{\partial n}\,\mathrm{d}S,
  \label{eq:gf-dirichlet}
\end{equation}
其中 $\phi=u|_{\Sigma}$.

罗宾条件 $\bigl[\alpha\partial u/\partial n+\beta u\bigr]_{\Sigma}=\psi$,并且 $G$ 满足同一系数的齐次条件.边界上 $\partial G/\partial n=-(\beta/\alpha)G$,$\partial u/\partial n=(\psi-\beta u)/\alpha$.代入\eqref{eq:gf-fundamental}后,含 $u$ 的面积分两两抵消,
\begin{equation}
  u(\br_0)
  =\iiint_{T} G f\,\mathrm{d}V
  -\frac{1}{\alpha}\iint_{\Sigma} G(\br,\br_0)\,\psi(\br)\,\mathrm{d}S.
  \label{eq:gf-robin}
\end{equation}

诺伊曼条件不能要求 $\nabla^2 G=\delta$ 并且 $\partial G/\partial n|_{\Sigma}=0$.对方程积分,左边是边界上法向导数的积分,等于零,右边是 $1$.物理上,绝热区域内放一个净热源,温度不会达到稳态.三维里改用
\begin{equation}
  \nabla^2 G=\delta(\br-\br_0)-\frac{1}{V_T},
  \qquad
  \frac{\partial G}{\partial n}\Big|_{\Sigma}=0,
\end{equation}
$V_T$ 是 $T$ 的体积.均匀汇把点源的总量抵消.二维用面积 $A_T$ 代替 $V_T$.原问题本身也要满足相容条件 $\iiint f\,\mathrm{d}V=\iint\partial u/\partial n\,\mathrm{d}S$.

\subsection{互易性}
\begin{equation}
  G(\br,\br_0)=G(\br_0,\br).
  \label{eq:gf-reciprocity}
\end{equation}
证明时把两个源点都挖成小球,对 $G(\cdot,\br_1)$ 和 $G(\cdot,\br_2)$ 用第二恒等式.体积内被积函数是零,外边界上齐次条件使面积分为零,两个小球的极限留下 $G(\br_1,\br_2)-G(\br_2,\br_1)=0$.

因此观察点和源点可以换写.\eqref{eq:gf-dirichlet}的常用形式是
\begin{equation}
  u(\br)
  =\iiint_{T} G(\br,\br_0)\,f(\br_0)\,\mathrm{d}V_0
  +\iint_{\Sigma}\phi(\br_0)\frac{\partial G}{\partial n_0}\,\mathrm{d}S_0.
  \label{eq:gf-dirichlet-obs}
\end{equation}
$f=0$ 时只剩边界积分:调和函数由边界值经 $\partial G/\partial n$ 再生.

\begin{example}
区间 $0<x<l$ 上 $u''=f$,$u(0)=u(l)=0$.
\end{example}
\begin{solution}
$G''=\delta(x-x_0)$,两端为零.两侧是直线,
\begin{equation}
  G=
  \begin{cases}
    A x, & 0\leq x\leq x_0, \\
    C(x-l), & x_0\leq x\leq l.
  \end{cases}
\end{equation}
连续:$A x_0=C(x_0-l)$.越过 $x_0$ 积分,$G'(x_0+)-G'(x_0-)=1$,即 $C-A=1$.解出
\begin{equation}
  G(x,x_0)=
  \begin{cases}
    (x_0-l)x/l, & x\leq x_0, \\
    x_0(x-l)/l, & x\geq x_0.
  \end{cases}
  \label{eq:gf-1d}
\end{equation}
它对称,而且在区间内部为负,和基本解 $-1/(4\pi R)$ 的符号一致.

取 $f=-1$,则 $u(x)=\int_0^l G(x,x_0)(-1)\,\mathrm{d}x_0=x(l-x)/2$.直接积分 $u''=-1$ 得到同一条抛物线.
\end{solution}

\section{电像法}
\idx{电像法}{dianxiangfa}
\idx{泊松核}{bosonghe}

\subsection{无界空间的基本解}
把格林函数拆成奇异部分和调和修正,
\begin{equation}
  G=G_0+G_1,
  \qquad
  \nabla^2 G_0=\delta(\br-\br_0),
  \qquad
  \nabla^2 G_1=0.
\end{equation}
狄利克雷问题还要求 $G_1|_{\Sigma}=-G_0|_{\Sigma}$.

三维里把源放在原点,$\nabla^2 G_0=0$（$r\neq 0$）的径向解是 $-C_1/r+C_2$.无穷远处趋于零则 $C_2=0$.对小球积分 $\nabla^2 G_0=\delta$,由散度定理
\begin{equation}
  \iint \frac{\partial G_0}{\partial r}\,\mathrm{d}S
  =4\pi C_1=1,
\end{equation}
所以 $C_1=1/(4\pi)$.源点移到 $\br_0$,
\begin{equation}
  G_0(\br,\br_0)=-\frac{1}{4\pi|\br-\br_0|}.
  \label{eq:gf-free-3d}
\end{equation}
二维同样得到
\begin{equation}
  G_0(\br,\br_0)=\frac{1}{2\pi}\ln|\br-\br_0|.
  \label{eq:gf-free-2d}
\end{equation}
静电里它们是自由空间中电荷 $-\varepsilon_0$ 的电位,方程是 $\nabla^2 V=-\rho/\varepsilon_0$ 时电荷要改号.

边界足够简单时,$G_1$ 可以写成区域外面若干像源的基本解.像源不在 $T$ 内,所以 $T$ 里仍然只有原来的点源,边界条件靠像源去配.

\subsection{接地球内}
球半径 $a$,源在距球心 $r_0<a$ 处,要求球面 $G=0$.像放在同一射线上的反演点
\begin{equation}
  \br_1=\frac{a^2}{r_0^2}\br_0,
  \qquad
  r_1=\frac{a^2}{r_0}>a.
\end{equation}
球面上 $|\br|=a$ 时,余弦定理给出 $|\br-\br_1|=(a/r_0)\,|\br-\br_0|$.像的强度取 $a/r_0$,两项相消,
\begin{equation}
  G(\br,\br_0)
  =-\frac{1}{4\pi|\br-\br_0|}
   +\frac{a/r_0}{4\pi\bigl|\br-(a^2/r_0^2)\br_0\bigr|}.
  \label{eq:gf-sphere}
\end{equation}
像在球外,球内 $\nabla^2 G$ 只有原来的 $\delta$.

\subsection{圆盘}
$\rho<a$,二维狄利克雷格林函数用同一个反演点 $\br_1$,
\begin{equation}
  G
  =\frac{1}{2\pi}\ln|\br-\br_0|
   -\frac{1}{2\pi}\ln|\br-\br_1|
   +\frac{1}{2\pi}\ln\frac{a}{r_0}.
  \label{eq:gf-disk}
\end{equation}
常数项不破坏调和性.圆周上 $|\br-\br_1|=(a/r_0)\,|\br-\br_0|$,对数相加是零.

\subsection{泊松积分}
$f=0$ 时,\eqref{eq:gf-dirichlet-obs}只剩 $\partial G/\partial n_0$.下面三个区域都可以把这个法向导数算出来.

\begin{example}
球内调和,$u|_{r=a}=f(\theta,\phi)$.
\end{example}
\begin{solution}
用\eqref{eq:gf-sphere}.记 $R=|\br-\br_0|$,$R_1=|\br-\br_1|$,夹角为 $\Theta$,$r_1=a^2/r_0$.则
\begin{equation}
  G=-\frac{1}{4\pi}\Biggl(\frac{1}{R}-\frac{a}{r_0 R_1}\Biggr).
\end{equation}
外法向是 $\partial/\partial r_0$.在 $r_0=a$ 上像点和源点重合,$R_1=R$,但是法向导数不是零.直接求导再令 $r_0=a$,用
\begin{equation}
  R^2=r^2+a^2-2ar\cos\Theta
\end{equation}
把 $a-r\cos\Theta=(R^2+a^2-r^2)/(2a)$ 代回去,两项里的 $1/(aR)$ 相消,留下
\begin{equation}
  \frac{\partial G}{\partial r_0}\Big|_{r_0=a}
  =\frac{1}{4\pi a}\frac{a^2-r^2}{R^3}.
\end{equation}
面积元 $a^2\sin\theta_0\,\mathrm{d}\theta_0\,\mathrm{d}\phi_0$,所以
\begin{equation}
  u(r,\theta,\phi)
  =\frac{a}{4\pi}\int_0^{\pi}\int_0^{2\pi}
    f(\theta_0,\phi_0)
    \frac{a^2-r^2}{(a^2-2ar\cos\Theta+r^2)^{3/2}}
    \sin\theta_0\,\mathrm{d}\theta_0\,\mathrm{d}\phi_0.
  \label{eq:gf-poisson-sphere}
\end{equation}
这是球的泊松积分.$r\to a$ 时核变成球面上的 $\delta$,$u$ 回到 $f$.球心处 $\cos\Theta$ 消失,$u(0)$ 等于 $f$ 在球面上的平均.
\end{solution}

\begin{example}
上半空间 $z>0$,$u|_{z=0}=f(x,y)$,无穷远处 $u\to 0$.
\end{example}
\begin{solution}
像电荷与源反号,放在 $(x_0,y_0,-z_0)$,
\begin{equation}
  G=-\frac{1}{4\pi|\br-\br_0|}+\frac{1}{4\pi|\br-\br_1|}.
\end{equation}
$z=0$ 上两个距离相等,$G=0$.上半空间的外法向是 $-\mathbf{e}_z$,所以 $\partial/\partial n_0=-\partial/\partial z_0$.在平面上
\begin{equation}
  \frac{\partial G}{\partial n_0}
  =\frac{z}{2\pi\bigl[(x-x_0)^2+(y-y_0)^2+z^2\bigr]^{3/2}}.
\end{equation}
因此
\begin{equation}
  u(x,y,z)
  =\frac{z}{2\pi}\iint_{\mathbb{R}^2}
    \frac{f(x_0,y_0)}
         {\bigl[(x-x_0)^2+(y-y_0)^2+z^2\bigr]^{3/2}}
    \,\mathrm{d}x_0\,\mathrm{d}y_0.
  \label{eq:gf-poisson-halfspace}
\end{equation}
\end{solution}

\begin{example}
圆盘 $\rho<a$ 与上半平面 $y>0$.
\end{example}
\begin{solution}
圆盘上对\eqref{eq:gf-disk}求外法向 $\partial/\partial\rho_0$,在 $\rho_0=a$ 处得到
\begin{equation}
  u(\rho,\phi)
  =\frac{a^2-\rho^2}{2\pi}\int_0^{2\pi}
    \frac{f(\phi_0)}
         {a^2-2a\rho\cos(\phi-\phi_0)+\rho^2}
    \,\mathrm{d}\phi_0.
  \label{eq:gf-poisson-disk}
\end{equation}
上半平面用反号像点 $(x_0,-y_0)$,外法向 $-\mathbf{e}_y$,
\begin{equation}
  u(x,y)
  =\frac{y}{\pi}\int_{-\infty}^{\infty}
    \frac{f(x_0)}{(x-x_0)^2+y^2}\,\mathrm{d}x_0.
  \label{eq:gf-poisson-halfplane}
\end{equation}
$f=1$ 时两个积分都等于 $1$,常值函数是调和的,边界值也是 $1$.
\end{solution}

\begin{example}
圆盘边界 $u(a,\phi)=A\cos\phi$.
\end{example}
\begin{solution}
把\eqref{eq:gf-poisson-disk}的核按 $\cos(n(\phi-\phi_0))$ 展开,或者直接用分离变量.内部有限的解是 $\rho^{n}$ 的傅里叶级数,边界只有 $\cos\phi$,
\begin{equation}
  u(\rho,\phi)=A\frac{\rho}{a}\cos\phi.
\end{equation}
它和柱坐标一章的圆内边值是同一个函数.泊松积分把分离变量的结果包进一个核对边界值的积分里.
\end{solution}

\section{含时问题的格林函数}

\subsection{波动方程}
受迫振动
\begin{equation}
  u_{tt}-a^2\nabla^2 u=f(\br,t),
\end{equation}
边界是\eqref{eq:gf-bc}那种线性条件,初值 $u|_{t=0}=\varphi$, $u_t|_{t=0}=\psi$.连续力是一串冲量,
\begin{equation}
  f(\br,t)
  =\iiint_{T}\int_{-\infty}^{\infty}
    f(\br_0,\tau)\,\delta(\br-\br_0)\,\delta(t-\tau)
    \,\mathrm{d}V_0\,\mathrm{d}\tau.
\end{equation}
格林函数是 $t_0>0$ 时单位冲量的响应,
\begin{equation}
  G_{tt}-a^2\nabla^2 G=\delta(\br-\br_0)\delta(t-t_0),
\end{equation}
边界条件齐次,并且 $G|_{t=0}=G_t|_{t=0}=0$.因果性要求 $t<t_0$ 时 $G=0$.

把原方程和 $G$ 的方程相减,在 $T\times[0,t]$ 上积分.时间部分用
\begin{equation}
  G u_{t_0 t_0}-u G_{t_0 t_0}
  =\partial_{t_0}\bigl(G u_{t_0}-u G_{t_0}\bigr),
\end{equation}
空间部分用格林第二恒等式,得到
\begin{equation}
\begin{aligned}
  u(\br,t)
  &=\iiint_{T}\int_0^{t} G\,f\,\mathrm{d}V_0\,\mathrm{d}t_0 \\
  &\quad+a^2\iint_{\Sigma}\int_0^{t}
      \Bigl(G\frac{\partial u}{\partial n_0}-u\frac{\partial G}{\partial n_0}\Bigr)
      \,\mathrm{d}S_0\,\mathrm{d}t_0 \\
  &\quad+\iiint_{T}\bigl[G u_{t_0}-u G_{t_0}\bigr]_{t_0=0}\,\mathrm{d}V_0.
\end{aligned}
\label{eq:gf-wave-rep}
\end{equation}
三项依次是外力、边界数据和初值.把 $G$ 的边界条件取得和问题同一类型,未知的边界组合就不再出现.

波动格林函数的互易性带时间反演,
\begin{equation}
  G(\br,t;\br_0,t_0)=G(\br_0,-t_0;\br,-t).
\end{equation}
原因是波动方程对 $t\to -t$ 不变,而因果性把未来的源变成过去的响应.

\subsection{热传导方程}
\begin{equation}
  u_t-a^2\nabla^2 u=f,
  \qquad
  u|_{t=0}=\varphi.
\end{equation}
热方程对时间是一阶的,同样的手续只留下一个初值项,
\begin{equation}
\begin{aligned}
  u(\br,t)
  &=\iiint_{T}\int_0^{t} G f\,\mathrm{d}V_0\,\mathrm{d}t_0 \\
  &\quad+a^2\iint_{\Sigma}\int_0^{t}
      \Bigl(G\frac{\partial u}{\partial n_0}-u\frac{\partial G}{\partial n_0}\Bigr)
      \,\mathrm{d}S_0\,\mathrm{d}t_0 \\
  &\quad+\iiint_{T}\bigl[u G\bigr]_{t_0=0}\,\mathrm{d}V_0.
\end{aligned}
\label{eq:gf-heat-rep}
\end{equation}

\subsection{冲量定理}
冲量定理把 $\delta(t-\tau)$ 的外力换成 $t=\tau$ 处初值的跳跃,然后去解齐次方程.对波动方程,零初位移和单位冲量意味着
\begin{equation}
  G\big|_{t=\tau+0}=0,
  \qquad
  G_t\big|_{t=\tau+0}=\delta(x-\xi).
\end{equation}
对热方程,跳跃发生在函数值本身,$G|_{t=\tau+0}=\delta(x-\xi)$.

\begin{example}
无限长弦,$u_{tt}-a^2 u_{xx}=f$,$u|_{t=0}=u_t|_{t=0}=0$.
\end{example}
\begin{solution}
达朗贝尔公式把时间原点挪到 $\tau$.初位移为零、初速度为 $\delta(x-\xi)$ 时,依赖区间里的积分是 $1/(2a)$,区间外是零,
\begin{equation}
  G(x,t;\xi,\tau)
  =
  \begin{cases}
    1/(2a), & |x-\xi|<a(t-\tau), \\
    0, & |x-\xi|>a(t-\tau).
  \end{cases}
\end{equation}
叠加
\begin{equation}
  u(x,t)
  =\int_0^{t}\int_{x-a(t-\tau)}^{x+a(t-\tau)}
    \frac{1}{2a}\,f(\xi,\tau)\,\mathrm{d}\xi\,\mathrm{d}\tau.
\end{equation}
这就是前面齐次化原理里的公式.扰动仍然不超过速度 $a$.
\end{solution}

\begin{example}
弦长 $\ell$,两端自由,$u_x|_{x=0}=u_x|_{x=\ell}=0$,初值为零,外力 $A\cos(\pi x/\ell)\sin\omega t$.
\end{example}
\begin{solution}
外力已经是 $n=1$ 的余弦模态.令
\begin{equation}
  u(x,t)=q(t)\cos\frac{\pi x}{\ell},
  \qquad
  \omega_0=\frac{\pi a}{\ell}.
\end{equation}
代入波动方程,
\begin{equation}
  q''+\omega_0^2 q=A\sin\omega t,
  \qquad
  q(0)=q'(0)=0.
\end{equation}
$\omega\neq\omega_0$ 时
\begin{equation}
  q(t)
  =\frac{A}{\omega_0^2-\omega^2}
    \Bigl(\sin\omega t-\frac{\omega}{\omega_0}\sin\omega_0 t\Bigr).
\end{equation}
$\omega=\omega_0$ 时分母为零,改求长期项.特解试 $t\cos\omega_0 t$,配上齐次解以满足初值,
\begin{equation}
  q(t)
  =\frac{A}{2\omega_0^2}
    \bigl(\sin\omega_0 t-\omega_0 t\cos\omega_0 t\bigr).
\end{equation}
振幅随 $t$ 线性增长,这是共振.同一结果也可以把诺伊曼余弦族的冲量格林函数
\begin{equation}
  G=\frac{t-\tau}{\ell}
  +\frac{2}{\pi a}\sum_{n=1}^{\infty}\frac{1}{n}
    \sin\frac{n\pi a(t-\tau)}{\ell}
    \cos\frac{n\pi\xi}{\ell}\cos\frac{n\pi x}{\ell}
\end{equation}
对 $f$ 积分,正交性只留下 $n=1$.
\end{solution}

\subsection{热核}
\idx{热核}{rehe}
无限长杆 $u_t-a^2 u_{xx}=f$,$u|_{t=0}=0$.冲量问题是齐次热方程,初值 $\delta(x-\xi)$.解是高斯核
\begin{equation}
  G(x,t;\xi,\tau)
  =\frac{1}{2a\sqrt{\pi(t-\tau)}}
    \exp\Bigl(-\frac{(x-\xi)^2}{4a^2(t-\tau)}\Bigr),
  \qquad t>\tau.
  \label{eq:gf-heat-kernel}
\end{equation}
因此
\begin{equation}
  u(x,t)
  =\int_0^{t}\int_{-\infty}^{\infty}
    f(\xi,\tau)\,G(x,t;\xi,\tau)\,\mathrm{d}\xi\,\mathrm{d}\tau.
\end{equation}
只要 $t>\tau$,$G$ 在整条直线上都不为零:信号的传播速度是无限的.这是傅里叶热传导定律的后果.时间间隔不是极短时,它仍然是好的近似.

半直线 $x>0$ 用延拓.端点 $u|_{x=0}=0$ 时把 $f$ 奇延拓,热核相减,
\begin{equation}
\begin{aligned}
  G
  &=\frac{1}{2a\sqrt{\pi(t-\tau)}}
    \Biggl[
      \exp\Bigl(-\frac{(x-\xi)^2}{4a^2(t-\tau)}\Bigr) \\
      &\qquad
      -\exp\Bigl(-\frac{(x+\xi)^2}{4a^2(t-\tau)}\Bigr)
    \Biggr].
\end{aligned}
\end{equation}
$x=0$ 时两项相同,$G=0$.端点绝热 $u_x|_{x=0}=0$ 时改成偶延拓,上式的减号换成加号.偶函数的导数在 $x=0$ 为零.

有限杆 $0<x<\ell$、两端为零时,把 $\delta(x-\xi)$ 展成正弦级数,每个模态再乘衰减因子,
\begin{equation}
  G
  =\frac{2}{\ell}\sum_{n=1}^{\infty}
    \sin\frac{n\pi x}{\ell}\sin\frac{n\pi\xi}{\ell}
    \exp\Bigl(-a^2\bigl(\tfrac{n\pi}{\ell}\bigr)^2(t-\tau)\Bigr).
\end{equation}
$t\to\tau^+$ 时它回到 $\delta(x-\xi)$.这和分离变量写出的解是同一件事.

\section{推广的格林公式}

前面的恒等式针对拉普拉斯算子.一般的二阶线性算子也有伴随算子,使 $v\,Lu-u\,Mv$ 成为散度.下面在两个自变量上写出这个结构,更多自变量的做法相同.

\subsection{伴随算子}
\begin{equation}
  L=a_{11}\partial_{xx}+2a_{12}\partial_{xy}+a_{22}\partial_{yy}
    +b_1\partial_x+b_2\partial_y+c.
\end{equation}
伴随算子是
\begin{equation}
\begin{aligned}
  Mv
  &=\partial_{xx}(a_{11}v)+2\partial_{xy}(a_{12}v)+\partial_{yy}(a_{22}v) \\
  &\quad-\partial_x(b_1 v)-\partial_y(b_2 v)+c v.
\end{aligned}
\end{equation}
$L=M$ 时称 $L$ 自伴.拉普拉斯算子、波动算子 $\partial_{tt}-a^2\nabla^2$ 都是自伴的.在平面区域 $T$ 上,
\begin{equation}
  \iint_{T}\bigl(v\,Lu-u\,Mv\bigr)\,\mathrm{d}S
  =\int_{\Sigma}\bigl[X\cos(n,x)+Y\cos(n,y)\bigr]\,\mathrm{d}s.
  \label{eq:gf-general}
\end{equation}
$X,Y$ 由分部积分定出,边界上只出现 $u,v$ 及其一阶导数.挖去源点处的小圆,令半径趋于零,就得到和\eqref{eq:gf-fundamental}同一形状的表示.差别只在边界积分里多出低阶项 $b_1,b_2$.

\subsection{椭圆型}
\begin{equation}
  Lu=u_{xx}+u_{yy}+b_1 u_x+b_2 u_y+c u=f.
\end{equation}
伴随是 $Mv=v_{xx}+v_{yy}-(b_1 v)_x-(b_2 v)_y+c v$.取 $MG=\delta(\br-\br_0)$,则
\begin{equation}
\begin{aligned}
  u(\br_0)
  &=\iint_{T} G f\,\mathrm{d}S \\
  &\quad-\int_{\Sigma}\Bigl[
      G\frac{\partial u}{\partial n}-u\frac{\partial G}{\partial n}
      +(b_1\cos(n,x)+b_2\cos(n,y))\,u G
    \Bigr]\,\mathrm{d}s.
\end{aligned}
\end{equation}
第一类边界让 $G|_{\Sigma}=0$,面积分只留下 $u\,\partial G/\partial n$.第二类边界让
\begin{equation}
  \Bigl[\frac{\partial G}{\partial n}
    -b_1\cos(n,x)\,G-b_2\cos(n,y)\,G\Bigr]_{\Sigma}=0,
\end{equation}
面积分只留下 $-G\,\partial u/\partial n$.$b_1=b_2=0$ 时回到泊松方程.

\begin{example}
单色波的基尔霍夫公式.
\end{example}
\begin{solution}
由 $U_{tt}-a^2\nabla^2 U=-f(\br)e^{\imath\omega t}$,令 $U=u e^{\imath\omega t}$,得到亥姆霍兹方程
\begin{equation}
  \nabla^2 u+k^2 u=f,
  \qquad k=\omega/a.
\end{equation}
它自伴.自由空间的一个格林函数是出射波
\begin{equation}
  G=\frac{1}{4\pi}\frac{e^{\imath k R}}{R},
  \qquad R=|\br-\br_0|.
\end{equation}
注意这里方程是 $(\nabla^2+k^2)u=f$,基本解的符号和 $\nabla^2( -1/(4\pi R) )=\delta$ 在 $k=0$ 时差一个负号: $k\to 0$ 时 $G\to +1/(4\pi R)$,对应的是 $\nabla^2 G=-\delta$.把这个 $G$ 代入第二恒等式,
\begin{equation}
\begin{aligned}
  u(\br)
  &=\frac{1}{4\pi}\iint_{\Sigma}\Biggl[
      \frac{e^{\imath k R}}{R}\frac{\partial u}{\partial n}
      -u\frac{\partial}{\partial n}\Bigl(\frac{e^{\imath k R}}{R}\Bigr)
    \Biggr]\,\mathrm{d}S \\
  &\quad-\frac{1}{4\pi}\iiint_{T}\frac{e^{\imath k R}}{R}f\,\mathrm{d}V.
\end{aligned}
\label{eq:gf-kirchhoff}
\end{equation}
这是基尔霍夫公式,衍射理论里单色波的积分表示.无源、无穷远辐射条件合适时,体积分消失,边界上的 $u$ 和 $\partial u/\partial n$ 决定整个场.
\end{solution}

\subsection{抛物型}
令 $y=a^2 t$,一维热算子写成 $Lu=u_{xx}-u_y$.伴随是 $M=\partial_{xx}+\partial_y$,时间导数反向.在矩形 $x\in[x_1,x_2]$, $0\leq y\leq a^2 t_0$ 上,格林公式的底边贡献初值,两条竖边贡献边界,
\begin{equation}
\begin{aligned}
  u(x_0,t_0)
  &=\int_{x_1}^{x_2}\int_0^{a^2 t_0} G f\,\mathrm{d}x\,\mathrm{d}y \\
  &\quad-\int_0^{a^2 t_0}
      \Bigl(G\frac{\partial u}{\partial x}-u\frac{\partial G}{\partial x}\Bigr)
      \Big|_{x_1}^{x_2}\,\mathrm{d}y \\
  &\quad-\int_{x_1}^{x_2}(u G)\big|_{y=0}\,\mathrm{d}x.
\end{aligned}
\end{equation}
两端为零时要求 $G|_{x_i}=0$;两端绝热时要求 $G_x|_{x_i}=0$.这就是\eqref{eq:gf-heat-rep}在一维上的样子.

\subsection{双曲型与黎曼函数}
\idx{黎曼函数}{limanhanshu}
\begin{equation}
  Lu=u_{xx}-u_{yy}+b_1 u_x+b_2 u_y+c u=f.
\end{equation}
伴随把一阶项改成 $-(b_i v)$ 的导数.柯西数据 $u$ 和 $\partial u/\partial n$ 给在曲线 $l$ 上.从观察点 $M_0(x_0,y_0)$ 作两条特征线 $x\pm y=\textrm{常数}$,交 $l$ 于 $M_1,M_2$,在特征三角形 $M_0M_1M_2$ 上积分.

黎曼函数 $v(x,y;x_0,y_0)$ 满足 $Mv=0$,$v(M_0)=1$,并且沿两条特征线
\begin{align}
  M_2M_0&:\quad
    \frac{\partial v}{\partial s}+\frac{b_2+b_1}{2\sqrt{2}}v=0, \\
  M_0M_1&:\quad
    \frac{\partial v}{\partial s}-\frac{b_2-b_1}{2\sqrt{2}}v=0.
\end{align}
这样选了以后,特征线上的积分只留下端点,得到
\begin{equation}
\begin{aligned}
  u(x_0,y_0)
  &=\half\bigl[u(M_1)v(M_1)+u(M_2)v(M_2)\bigr] \\
  &\quad+\half\int_{M_1M_2}
      \bigl[X\cos(n,x)+Y\cos(n,y)\bigr]\,\mathrm{d}s \\
  &\quad-\half\iint_{M_0M_1M_2} v\,f\,\mathrm{d}S.
\end{aligned}
\label{eq:gf-riemann}
\end{equation}
黎曼函数对称,$v(x,y;x_0,y_0)=v(x_0,y_0;x,y)$.若 $l$ 上的柯西数据为零,并且 $f=-2\delta(x-\xi)\delta(y-\eta)$,则 $u(x_0,y_0)=v(\xi,\eta;x_0,y_0)$.所以黎曼函数就是这个双曲算子的点源影响函数.

\begin{example}
常系数方程 $u_{xx}-u_{yy}=0$ 的黎曼函数,以及它和达朗贝尔公式的关系.
\end{example}
\begin{solution}
此时 $b_1=b_2=c=0$,伴随和 $L$ 相同.特征线上的传输条件变成 $\partial v/\partial s=0$,再加 $v(M_0)=1$,所以 $v\equiv 1$.

把 $y$ 看成时间,$x$ 看成空间,特征线斜率是 $\pm 1$,对应波速 $a=1$.初值给在 $y=0$:$u(x,0)=\varphi(x)$, $u_y(x,0)=\psi(x)$.过 $(x_0,y_0)$ 的两条特征线交 $y=0$ 于 $x_0-y_0$ 和 $x_0+y_0$.\eqref{eq:gf-riemann}里 $v=1$,$f=0$,底边上的 $X,Y$ 化为初值的积分,结果是
\begin{equation}
  u(x_0,y_0)
  =\frac{\varphi(x_0-y_0)+\varphi(x_0+y_0)}{2}
  +\half\int_{x_0-y_0}^{x_0+y_0}\psi(s)\,\mathrm{d}s.
\end{equation}
这就是 $a=1$ 的达朗贝尔公式.波速为 $a$ 时,把 $y$ 换成 $at$,积分前面的 $1/2$ 变成 $1/(2a)$.
\end{solution}

本章泊松方程是 $\nabla^2 u=f$,自由空间基本解是 $-1/(4\pi R)$.若题目里的方程改成 $-\nabla^2 u=\rho$,格林函数要整体变号.二维基本解是 $(1/(2\pi))\ln R$.

\begin{problemset}
\renewcommand{\theenumi}{\thechapter.\arabic{enumi}}
\renewcommand{\labelenumi}{\theenumi.}
\item 用镜像法求圆盘 $\rho<a$ 上的狄利克雷格林函数,边界值 $u|_{\rho=a}=f(\phi)$.由此得到泊松积分.

\item 用镜像法解半平面 $y>0$ 上的 $\nabla^2 u=0$, $u(x,0)=f(x)$.

\item 圆盘 $\rho\leq a$ 上的调和函数,边界分别为 $u|_{\rho=a}=A\cos\phi$ 和 $u|_{\rho=a}=A+B\sin\phi$.分别求盘内的 $u$.

\item 把圆盘泊松核展成傅里叶级数,逐项积分,并与分离变量的结果比较.

\item 区间 $0<x<\ell$ 上 $u''=f$, $u(0)=u(\ell)=0$.求格林函数,并验证它在 $x=x_0$ 处连续、导数跳跃为 $1$,而且在区间内为负.

\item 两端固定的弦受 $\rho f=\rho\Phi(x)\sin\omega t$ 驱动.写出受迫解,讨论共振,并在共振时写出解.

\item 同一根弦在内点 $x_0$ 受到集中谐力.求位移.

\item 解 $u_t-a^2 u_{xx}=0$, $x>0$, $u(0,t)=f(t)$, $u(x,0)=\varphi(x)$.说明狄利克雷端点为什么用奇延拓.

\item 解 $u_t-a^2 u_{xx}=0$, $x>0$, $u_x(0,t)=q(t)$, $u(x,0)=0$.

\item 用推广的格林公式写出
\begin{equation*}
  u_{tt}-a^2 u_{xx}=f(x,t),
  \qquad
  u(x,0)=\varphi(x),
  \qquad
  u_t(x,0)=\psi(x)
\end{equation*}
的解,并在 $f=0$ 时核对达朗贝尔公式.

\item 对 $u_{xx}-u_{yy}=0$ 取黎曼函数 $v\equiv 1$,说明它如何回到达朗贝尔公式.
\end{problemset}

